%%%%%%%%%%%%%%%%%%%%%%%%%%%%%%%%%%%%%%%%%%%%%%%%%%%%%%%%%%%%%%%%%%%%%%%%%%%
% Lehreinheit 7: Kognitive Verzerrungen
% Quelle: WuZ-Buch, Kap. 15 (LE 7), Kap. 10 (Charpentier-Gewichtsillusion),
%         Kap. 4 (Confirmation Bias, HCEG, Misdirection), Kap. 5 (Alief),
%         Kap. 7 (Noise und Entscheidungshygiene), Kap. 5b (Congruence Bias)
% Staunmoment: Charpentier-Illusion  |  Kernbegriff: Bias vs. Noise
% Die Charpentier-Illusion wird im Buch vollstaendig erklaert und darf
% hier beschrieben werden.
%%%%%%%%%%%%%%%%%%%%%%%%%%%%%%%%%%%%%%%%%%%%%%%%%%%%%%%%%%%%%%%%%%%%%%%%%%%
\nsl{LE~5 hat gezeigt, dass wir lieber bestätigen als prüfen, LE~6, dass
ohne Vergleich selbst grobe Täuschungen unentdeckt bleiben. Diese Einheit
fragt, woher die Neigung zur Täuschung kommt und wie sie in die Forschung
gelangt. Sie beginnt mit drei Streichholzschachteln, die jede Person
selbst in die Hand nimmt. Die Täuschung, die dabei entsteht, verschwindet
auch dann nicht, wenn man ihre Ursache kennt. Daraus entwickelt die
Einheit ihre Leitthese: Verzerrung ist der Normalbetrieb eines
vorhersagenden Gehirns, kein Zeichen von Dummheit. Bändigen lässt sie
sich deshalb nur durch Methode, nicht durch guten Willen. Die Einheit
ordnet die wichtigsten Verzerrungen den Phasen eines Forschungsprozesses
zu, behandelt Bestätigungsfehler, Hot-Cold-Empathy-Gap und Misdirection,
unterscheidet Bias von Noise und stellt die methodischen Gegenmittel vor.
Die Replikationskrise zeigt, was geschieht, wenn eine ganze Disziplin
diese Gegenmittel zu wenig nutzt.
\newpage}
\os{\begin{description}
\item[{\rot\bf Heute}]~\\[-\lblskip]
\bi
  \ite drei Streichholzschachteln, die sich nicht an die Physik halten
  \ite Bias ist Normalbetrieb: Heuristiken und Vorhersagen
  \ite Verzerrungen im Forschungsprozess: von der Planung bis zur
    Publikation
  \ite Bias und Noise -- und was dagegen hilft
  \ite die Replikationskrise
  \ite Bias-Jagd in einem echten Artikel
\ei
\end{description}
\newpage}

\begin{vorbereitung}[Charpentier-Illusion]
\textbf{Eine Woche vorher}
\begin{checkliste}
  \item Sets bauen: je Set drei gleiche Streichholzschachteln derselben
    Marke, Hölzer entfernen. In die \emph{oberste} Schachtel zwei
    AA-Batterien legen und mit etwas Papier oder Watte gegen Klappern
    sichern, sodass sie zusammen ca.\,50\,g wiegt. Die beiden anderen
    bleiben leer (je ca.\,3\,g). Die Schachteln mit einem Gummiband zum
    Stapel binden. Die beschwerte Schachtel unauffällig markieren (Punkt
    innen an der Lade), damit beim Zusammenbauen nichts vertauscht wird
    (\buch{Kap.~10, Charpentier-Gewichtsillusion}).
  \item Ein Set pro vier bis fünf Personen, plus zwei Ersatzsets; bei
    30~Teilnehmenden also acht bis neun Sets.
  \item Eine Küchenwaage (1-g-Auflösung) für die Auflösung einpacken.
  \item Den Ablauf an einer Testperson ausprobieren: Erst den Stapel heben,
    dann das Gewicht in Gramm schätzen und notieren, dann Gummiband
    lösen und nur die oberste Schachtel heben.
  \item \textbf{Artikel für die Bias-Jagd auswählen} (siehe unten) und
    eine Woche vorher als PDF verschicken mit der Bitte, ihn einmal zu
    lesen.
\end{checkliste}

\textbf{Auswahlkriterien für den Artikel der Bias-Jagd}
\begin{checkliste}
  \item empirische Originalarbeit (kein Review, keine Meta-Analyse), im
    Umfang von höchstens etwa zwölf Seiten;
  \item ohne Spezialwissen verständlich, möglichst aus dem Fach oder
    Interessengebiet der Gruppe; Englisch ist zumutbar;
  \item vollständiger Methodenteil mit Stichprobe, Zuteilung,
    Messinstrumenten und Auswertung, damit Bias-anfällige Stellen
    überhaupt erkennbar sind;
  \item mindestens zwei echte Angriffspunkte, z.\,B.\ kleine Stichprobe
    ohne Fallzahlplanung, viele abhängige Variablen, nicht blinde
    Beurteilung, fehlende Präregistrierung, Interpretation, die über die
    Daten hinausgeht;
  \item kein Artikel, der wegen Fälschung zurückgezogen wurde: Gesucht
    sind gewöhnliche Schwächen, keine Skandale;
  \item optional als Kontrast: ein präregistrierter Artikel oder ein
    \emph{Registered Report} zur selben Frage.
\end{checkliste}

\textbf{Handout und Übungsblatt}
\begin{checkliste}
  \item Handout \enquote{Kognitive Verzerrungen} sowie Übungsblatt LE~7
    Teil~A und Teil~B je einmal pro Person kopieren (Druckvorlage
    \texttt{handouts.pdf}, Seiten zu LE~7). Teil~A liegt zu Beginn aus
    und enthält einen fiktiven Kurztext zum Aufwärmen und das Raster für
    die Bias-Jagd. Teil~B (Selbstbezug, Tagebuch) wird erst zum
    Selbstbezug ausgeteilt, weil die Tagebuchfrage auf den Staunmoment
    anspielt.
  \item Den gewählten Artikel einmal pro Zweiergruppe ausdrucken, falls
    nicht alle ein Gerät dabeihaben.
  \item Für jede Person einen Zettel für die Gewichtsschätzung.
\end{checkliste}

\textbf{Am Tag selbst}
\begin{checkliste}
  \item Sets vor der Stunde auf die Tische verteilen, Stapel gebunden,
    beschwerte Schachtel oben.
  \item Tafel vorbereiten: links eine Strichliste \enquote{oberste
    Schachtel allein fühlt sich schwerer an: ja / nein}, rechts Platz für
    die notierten Schätzungen des Stapelgewichts.
\end{checkliste}

\textbf{Risiken und Plan B.} Wer die Illusion kennt, erlebt sie trotzdem.
Auch das gehört zur Lektion. Wer
nach dem Lösen des Gummibands versehentlich eine untere Schachtel hebt,
erlebt nichts Besonderes; deshalb den Ablauf vorher einmal ansagen. Fehlt
der Artikel oder hat die Gruppe ihn nicht gelesen, trägt der fiktive
Kurztext auf dem Übungsblatt (Teil~A) die Bias-Jagd in der Stunde, und der
Artikel wird zur Hausaufgabe.
\end{vorbereitung}

\begin{ablauf}
\begin{description}
\item[0--5 · Blitzlicht.] Tagebuch-Rückblick: \enquote{Wo haben Sie sich
  diese Woche selbst ertappt?}
\item[5--15 · Staunmoment.] Charpentier-Illusion durch Herumreichen;
  stille Gewichtsschätzung; Strichliste; Auflösung mit der Waage;
  zweiter Versuch mit Wissen. Staunen darüber, dass Wissen die Täuschung
  nicht auflöst (Alief, LE~1).
\item[15--50 · Theorie-Input.] Bias-Systematik; Confirmation Bias; HCEG;
  Misdirection; Noise und Entscheidungshygiene; Gegenmittel
  (Blind-Design, Präregistrierung); Replikationskrise.
\item[50--80 · Bias-Jagd.] Aufwärmen am fiktiven Kurztext (5~Min.), dann
  identifizieren Zweiergruppen im ausgewählten Artikel mindestens zwei
  Bias-anfällige Stellen und begründen sie (Raster auf dem Übungsblatt
  Teil~A, Aufgaben~1 und~2).
\item[80--90 · Selbstbezug.] Übungsblatt Teil~B austeilen. Welcher Bias
  droht im \emph{eigenen} geplanten Forschungsprozess, und wie begrenze
  ich ihn? (Aufgabe~3)
\end{description}
\end{ablauf}

%==============================================================================
\section{Staunmoment: Drei Schachteln}
%==============================================================================
\begin{staunbox}{Charpentier-Illusion}
Heben Sie den Stapel aus drei Schachteln an. \emph{Wie schwer ist er?
Schreiben Sie eine Zahl in Gramm auf.}

\medskip
Lösen Sie das Gummiband. Heben Sie \emph{nur die oberste} Schachtel.

\medskip
\emph{Was ist schwerer: die eine Schachtel oder alle drei?}
\end{staunbox}
\os{\newpage}

\begin{description}
\item[{\rot\bf Was ist passiert?}]~\\[-\lblskip]
\bi
  \ite Strichliste: Bei wem fühlte sich die einzelne Schachtel schwerer
    an als der Stapel?
  \ite Ihre Schätzungen des Stapelgewichts: Wie weit liegen sie
    auseinander?
  \ite Die Waage: Stapel ca.\ 56\,g, oberste Schachtel ca.\ 50\,g
  \ite Und jetzt, wo Sie es wissen: noch einmal heben
\ei
\os{\newpage}
\end{description}

\nsl{Die oberste Schachtel enthält zwei Batterien und wiegt etwa 50\,g,
die beiden anderen je etwa 3\,g. Der Stapel ist also größer \emph{und}
schwerer als die einzelne Schachtel, und trotzdem fühlt sich die einzelne
Schachtel schwerer an. Die Erklärung des Buchs (\buch{Kap.~10,
Charpentier-Gewichtsillusion}): Beim Heben einer einzelnen von drei
gleich aussehenden Schachteln erwartet das Gehirn etwa ein Drittel des
Gesamtgewichts, bekommt aber fast das ganze. Die enttäuschte Erwartung
erzeugt die Illusion. Augustin Charpentier hat diese Größe-Gewicht-Illusion
1891 beschrieben \parencite{charpentier1891}; sie gehört zu den
stabilsten Täuschungen der Wahrnehmungspsychologie. Eine einzige
akzeptierte Erklärung gibt es bis heute nicht, die meisten Ansätze teilen
aber den Kern, dass Erwartungen über das Gewicht die Wahrnehmung formen
\parencite{buckingham2014}.

Für diese Einheit ist ein zweiter Befund wichtiger. Flanagan und
Beltzner ließen Versuchspersonen gleich schwere Würfel unterschiedlicher
Größe immer wieder heben. Die Griffkräfte passten sich nach wenigen
Versuchen dem tatsächlichen Gewicht an, die Illusion blieb in voller
Stärke bestehen \parencite{flanagan2000independence}. Die Hand lernt, die
Wahrnehmung nicht. Das ist die Unterscheidung von Belief und Alief aus
LE~1 \parencite{gendler2008}: Sie \emph{wissen} jetzt, dass die Schachtel
leichter ist als der Stapel, und Sie \emph{spüren} trotzdem das Gegenteil.}

\begin{description}
\item[{\rot\bf Erwartung formt Wahrnehmung}]~\\[-\lblskip]
\begin{center}
\resizebox{0.82\textwidth}{!}{%
\begin{tikzpicture}[font=\sffamily,>=Stealth,
  box/.style={draw=rwuviolet,thick,fill=rwuviolet!8,rounded corners=2pt,
    minimum width=36mm,minimum height=11mm,align=center}]
  \node[box] (e) at (0,0) {Erwartung\\\enquote{ein Drittel von 56\,g}};
  \node[box] (s) at (8,0) {Hand meldet\\ca.\ 50\,g};
  \node[box,draw=rwucyan,fill=rwucyan40!40] (w) at (4,-2.5) {Wahrnehmung:\\\enquote{viel schwerer als gedacht}};
  \draw[->,thick] (e) -- (w);
  \draw[->,thick] (s) -- (w);
  \node[rwugray,align=center] at (4,0) {Abweichung\\groß};
\end{tikzpicture}}
\end{center}
\bi
  \ite die Hand passt sich an, die Wahrnehmung nicht (Flanagan \&
    Beltzner 2000)
  \ite Wissen löst die Täuschung nicht auf: Belief gegen Alief
\ei
\os{\newpage}
\end{description}

%==============================================================================
\section{Verzerrung als Normalbetrieb}
%==============================================================================
\nsl{Die Schachteln zeigen im Kleinen, was für das Urteilen allgemein
gilt. Das Gehirn arbeitet mit Vorhersagen und Faustregeln, weil es sonst
nie zu einem Urteil käme. Tversky und Kahneman haben 1974 gezeigt, dass
diese Heuristiken meist gut funktionieren, aber zu systematischen,
vorhersagbaren Fehlern führen, und zwar auch bei Fachleuten
\parencite{tversky1974judgment}. Kahneman fasst das später in das Bild
zweier Denkweisen: ein schnelles, automatisches System~1 und ein
langsames, anstrengendes System~2 \parencite{kahneman2011}. Verzerrungen
entstehen vor allem dort, wo System~1 antwortet und System~2 nicht
nachprüft. Eine Verzerrung ist damit kein Charakterfehler. Sie ist die
Rückseite einer Arbeitsweise, die im Alltag meistens trägt. Daraus folgt
die Leitthese der Einheit: Gegen Verzerrungen hilft ein Verfahren, das auch dann funktioniert,
wenn man voreingenommen ist; der Vorsatz, unvoreingenommen zu sein,
reicht nicht.}
\begin{description}
\item[{\rot\bf Bias}]~
\end{description}
\defn{bias}{Eine \textbf{kognitive Verzerrung} (Bias) ist eine
systematische, vorhersagbare Abweichung eines Urteils von einem
angemessenen Maßstab. Sie zeigt sich bei vielen Personen in dieselbe
Richtung.}
\bi
  \ite Heuristiken sparen Zeit und liegen meist richtig -- aber
    systematisch daneben, wo ihre Voraussetzung nicht stimmt
  \ite System~1 antwortet schnell, System~2 prüft (oder eben nicht)
  \ite auch Fachleute sind nicht ausgenommen
\ei
\os{\newpage}

\begin{merkbox}
Verzerrung ist kein Zeichen von Dummheit, sondern der Normalbetrieb eines
vorhersagenden Gehirns. Bändigen lässt sie sich nur durch Methode, nicht
durch guten Willen.
\end{merkbox}
\os{\newpage}

%==============================================================================
\section{Verzerrungen im Forschungsprozess}
%==============================================================================
\nsl{Für das wissenschaftliche Arbeiten ist eine Ordnung nach den Phasen
des Forschungsprozesses nützlicher als eine alphabetische Liste. In der
\emph{Planung} wirkt vor allem der Bestätigungsfehler: Man wählt Fragen,
Methoden und Literatur so, dass die eigene Vermutung gute Chancen hat
(LE~3 und 5). In der \emph{Erhebung} verschieben Erwartungen der
Forschenden die Daten: Studierende, denen man sagte, ihre Ratten seien
besonders labyrinthklug, maßen bei zufällig zugeteilten Tieren bessere
Leistungen \parencite{rosenthal1963experimenter}. Auch die erste Zahl in
einem Fragebogen kann als Anker alle weiteren Antworten prägen
\parencite{tversky1974judgment}. In der \emph{Auswertung} erlauben die
vielen kleinen Entscheidungen (welche Fälle ausschließen, welche Variable
kontrollieren, wann aufhören zu erheben) fast jedes gewünschte Ergebnis
\parencite{simmons2011falsepositive, gelmanloken2014forking}. In der
\emph{Interpretation} wirkt der Rückschaufehler: Wer das Ergebnis kennt,
hält es für vorhersehbar und glaubt, es schon immer erwartet zu haben
\parencite{fischhoff1975hindsight}; nachträglich formulierte Hypothesen
werden dann als vorab geplant ausgegeben \parencite{kerr1998harking}. In
der \emph{Publikation} schließlich verschwinden Nullbefunde in der
Schublade, und die veröffentlichte Literatur überschätzt Effekte
\parencite{rosenthal1979filedrawer}. LE~8 und LE~10 vertiefen die
statistische Seite.}
\begin{description}
\item[{\rot\bf Bias entlang des Forschungsprozesses}]~\\[-\lblskip]
\begin{center}
\small
\renewcommand{\arraystretch}{1.2}
\begin{tabular}{p{0.17\textwidth}p{0.3\textwidth}p{0.39\textwidth}}
\toprule
\textbf{Phase} & \textbf{Verzerrung} & \textbf{so sieht es aus} \\
\midrule
Planung & Bestätigungsfehler & nur Literatur und Tests, die passen \\
Erhebung & Versuchsleitereffekt, Anker & Erwartung färbt Messung; erste Zahl prägt \\
Auswertung & p-Hacking, Forking Paths & so lange analysieren, bis es passt \\
Interpretation & Rückschaufehler, HARKing & \enquote{Das war doch klar.} \\
Publikation & Publication Bias & Nullbefunde bleiben in der Schublade \\
\bottomrule
\end{tabular}
\end{center}
\os{\newpage}
\end{description}

%==============================================================================
\section{Der Bestätigungsfehler}
%==============================================================================
\nsl{Der Bestätigungsfehler ist die am besten untersuchte und vermutlich
folgenreichste Verzerrung \parencite{nickerson1998ubiquitous}. Das Buch
beschreibt ihn als Filter: Informationen werden so ausgewählt, gewichtet
und erinnert, dass sie die bestehende Meinung bestätigen
(\buch{Kap.~4, Bestätigungsfehler}). Er wirkt nicht nur beim Suchen
(der Congruence Bias aus LE~5), sondern auch beim Bewerten: Ein
unerwünschter Befund wird strenger geprüft als ein erwünschter. Das
historische Lehrstück ist Ignaz Semmelweis. Er senkte 1847 die
Müttersterblichkeit an der Ersten Geburtshilflichen Klinik in Wien
durch Händewaschen mit Chlorkalklösung von zeitweise über 15~Prozent auf
rund ein bis zwei Prozent \parencite{semmelweis1861}. Sein Befund
widersprach dem herrschenden Krankheitsverständnis und kränkte die
Standesehre der Ärzte; statt die Evidenz anzuerkennen, suchte man Gründe,
sie abzulehnen. Nicht das Fehlen von Daten war das Hindernis, sondern die
Abwehr unbequemer Daten.}
\begin{description}
\item[{\rot\bf Bestätigungsfehler}]~\\[-\lblskip]
\bi
  \ite \textbf{Suchen:} nur Belege, die passen (2-4-6-Test, LE~5)
  \ite \textbf{Bewerten:} unerwünschte Befunde werden strenger geprüft
  \ite \textbf{Erinnern:} Treffer bleiben, Fehlschläge verblassen
  \ite Semmelweis 1847: Händewaschen senkt die Müttersterblichkeit von
    zeitweise über 15\,\% auf rund 1--2\,\% -- und wird abgelehnt
\ei
\os{\newpage}
\end{description}

%==============================================================================
\section{Der Hot-Cold-Empathy-Gap}
%==============================================================================
\nsl{Eine zweite Quelle von Verzerrung ist der eigene Zustand. George
Loewenstein hat gezeigt, dass Menschen den Einfluss von Hunger, Schmerz,
Wut, Erschöpfung oder Erregung auf ihr eigenes Urteilen systematisch
unterschätzen \parencite{loewenstein1996visceral}. Wer ruhig ist (kalter
Zustand), kann nicht verlässlich simulieren, wie er im erregten
(heißen) Zustand entscheiden würde, und umgekehrt. Das Buch unterscheidet
drei Richtungen dieser Lücke und leitet daraus zwei Folgen für die
Forschung ab (\buch{Kap.~4, HCEG}): Erstens beschreiben Probanden, die im
ruhigen Labor ihr Verhalten in emotionalen Situationen vorhersagen, eher
ihr rationalisiertes Selbstbild als ihr tatsächliches Verhalten
(\emph{hypothetischer Bias}). Zweitens sind Forschende selbst nicht immer
im kalten Zustand; Publikationsdruck ist ein institutioneller heißer
Zustand, in dem Bestätigungsfehler und motiviertes Denken besonders
leicht greifen.}
\begin{description}
\item[{\rot\bf Hot-Cold-Empathy-Gap}]~\\[-\lblskip]
\bi
  \ite \textbf{kalt $\to$ heiß:} im ruhigen Zustand unterschätzt man
    künftige Impulse
  \ite \textbf{heiß $\to$ kalt:} im erregten Zustand hält man das
    Augenblicksziel für die stabile Präferenz
  \ite \textbf{zwischenmenschlich:} wer ruhig ist, unterschätzt den
    Zustand anderer
  \ite Forschung: hypothetischer Bias in Umfragen; Publikationsdruck als
    heißer Zustand der Forschenden
\ei
\os{\newpage}
\end{description}

%==============================================================================
\section{Misdirection}
%==============================================================================
\nsl{Die Zauberkunst nutzt dieselben Mechanismen gezielt. Taylor Hughes
beschreibt drei Werkzeuge der Misdirection, die auf der Bühne ebenso
wirken wie in Werbung und Politik \parencite{hughes2023misdirection}: die
Kontrolle des Gesprächs (wer die Fragen stellt, steuert das Denken), die
Verdichtung der Zeit (wer keine Zeit zum Nachdenken hat, findet keine
Alternativerklärung) und die Bestätigungsverstärkung (ein Publikum, das
eine Deutung angenommen hat, verteidigt sie und vollendet die Illusion
selbst). Das Buch liest Misdirection als gezielt falsch gefüttertes
Vorhersagesystem: Der Zauberkünstler setzt einen falschen Prior und sorgt
dafür, dass die Sinnesdaten ihn scheinbar bestätigen
(\buch{Kap.~4, Misdirection}). Für die Forschung ist das unbequem, denn
die Rolle des Zauberkünstlers übernimmt oft die eigene Fragestellung: Wie
eine Frage gerahmt ist, bestimmt, welche Antworten überhaupt in den Blick
kommen. Hughes' Gegenmittel ist schlicht: mehr Fragen stellen, und zwar
vor allem die, die gerade nicht gestellt werden.}
\begin{description}
\item[{\rot\bf Misdirection -- im Theater und in der Forschung}]~\\[-\lblskip]
\bi
  \ite \textbf{Gesprächskontrolle:} die erste Frage bestimmt den
    Denkraum -- auch die eigene Forschungsfrage
  \ite \textbf{Zeitdruck:} keine Zeit für Alternativerklärungen
  \ite \textbf{Bestätigungsverstärkung:} wer eine Deutung hat, vollendet
    die Illusion selbst
  \ite Gegenmittel: \emph{Welche Frage wird hier nicht gestellt? Welche
    Optionen sehe ich nicht?}
\ei
\os{\newpage}
\end{description}

%==============================================================================
\section{Bias und Noise}
%==============================================================================
\nsl{Nicht jeder Urteilsfehler ist eine Verzerrung. Kahneman, Sibony und
Sunstein unterscheiden \emph{Bias}, die systematische Abweichung in eine
Richtung, von \emph{Noise}, der unsystematischen Streuung von Urteilen,
die eigentlich gleich ausfallen sollten \parencite{kahneman2021}. Zwei
Gutachter bewerten dieselbe Arbeit unterschiedlich, dieselbe Person
bewertet sie an zwei Tagen unterschiedlich. Die Schachteln haben beides
gezeigt: Dass sich die einzelne Schachtel bei fast allen schwerer anfühlte,
ist Bias; dass die notierten Schätzungen des Stapelgewichts weit
auseinanderlagen, ist Noise. Gegen Noise empfehlen die Autoren
\emph{Entscheidungshygiene}: unabhängige Einschätzungen vor jeder
Diskussion einholen, strukturierte Bewertungsraster verwenden, mehrere
Urteile zusammenfassen, Informationssammlung und Bewertung trennen und
die Rolle des Advocatus Diaboli fest zuweisen, statt sie dem Zufall zu
überlassen \parencite{schulzhardt2002}. Für das wissenschaftliche Arbeiten
heißt das unter anderem: Kodierungen und Bewertungen von mindestens zwei
Personen unabhängig vornehmen lassen und ihre Übereinstimmung prüfen.}
\begin{description}
\item[{\rot\bf Bias und Noise}]~\\[-\lblskip]
\begin{center}
\resizebox{0.78\textwidth}{!}{%
\begin{tikzpicture}[font=\sffamily]
  \foreach \x/\t in {0/{\textbf{Bias}\\systematisch daneben},6.5/{\textbf{Noise}\\unsystematisch gestreut}}{
    \foreach \r in {1.6,1.1,0.6}{\draw[rwugray] (\x,0) circle (\r);}
    \fill[rwugray] (\x,0) circle (0.12);
    \node[below,align=center] at (\x,-1.8) {\t};
  }
  \foreach \p in {(0.85,0.75),(1.0,0.55),(0.75,0.5),(0.95,0.85),(1.1,0.7)}{
    \fill[rwuviolet] \p circle (0.11);}
  \foreach \p in {(6.5+0.9,0.2),(6.5-1.1,0.6),(6.5+0.3,-1.2),(6.5-0.5,-0.4),(6.5+0.6,1.1)}{
    \fill[rwucyan] \p circle (0.11);}
\end{tikzpicture}}
\end{center}
\bi
  \ite Schachteln: \enquote{einzelne schwerer} bei fast allen = Bias;
    weit gestreute Schätzungen = Noise
\ei
\os{\newpage}

\item[{\rot\bf Entscheidungshygiene}]~\\[-\lblskip]
\bi
  \ite unabhängige Einschätzungen \emph{vor} der Diskussion einsammeln
  \ite strukturierte Bewertungsraster statt Gesamteindruck
  \ite mehrere Urteile zusammenfassen
  \ite Informationssammlung und Bewertung trennen
  \ite Advocatus Diaboli als zugewiesene Rolle
  \ite in der Forschung: zwei unabhängige Kodierende, Übereinstimmung
    prüfen
\ei
\os{\newpage}
\end{description}

%==============================================================================
\section{Gegenmittel: Methode statt guter Wille}
%==============================================================================
\nsl{Weil Verzerrungen auch dann wirken, wenn man von ihnen weiß, setzen
die wirksamen Gegenmittel beim Verfahren an. \emph{Blind- und Doppelblind-Designs} nehmen Teilnehmenden und
Messenden das Wissen, das ihre Erwartungen in die Daten tragen würde; so
wurde etwa die Behauptung geprüft, Astrologen könnten Horoskope den
passenden Persönlichkeitsprofilen zuordnen
\parencite{carlson1985doubleblind}. Die \emph{Präregistrierung} legt
Hypothesen, Messgrößen und Auswertung fest, bevor die Daten erhoben sind,
und trennt damit sichtbar zwischen geprüften und nachträglich gefundenen
Ergebnissen \parencite{nosek2018preregistration}. In der
\emph{adversariellen Kollaboration} planen Forschende mit
gegensätzlichen Erwartungen eine Studie gemeinsam und einigen sich vorab,
welches Ergebnis welche Position stützen würde. \emph{Replikation} und
offen geteilte Daten machen die Prüfung unabhängig von denen, die das
Ergebnis erhofft haben. Die Feldstudie aus LE~6 ist ein kleines Beispiel
für den Wert dieser Disziplin: Ihr Zwischenstand zeigte die erhoffte
Abstufung, die mit wachsender Stichprobe verschwand, und das Buch hat den
eigenen frühen Trend ausdrücklich widerrufen (\buch{Kap.~12f}).}
\begin{description}
\item[{\rot\bf Welches Mittel gegen welchen Bias?}]~\\[-\lblskip]
\begin{center}
\small
\renewcommand{\arraystretch}{1.2}
\begin{tabular}{p{0.3\textwidth}p{0.55\textwidth}}
\toprule
\textbf{Verzerrung} & \textbf{Gegenmittel} \\
\midrule
Versuchsleitereffekt & Blind- und Doppelblind-Design \\
p-Hacking, HARKing & Präregistrierung \\
Bestätigungsfehler & Alternativhypothese, adversarielle Kollaboration \\
Noise in Bewertungen & unabhängige Mehrfachbeurteilung, Raster \\
Publication Bias & Registered Reports, Nullbefunde veröffentlichen \\
Zufallsbefund & Replikation, Daten teilen \\
\bottomrule
\end{tabular}
\end{center}
\os{\newpage}
\end{description}

%==============================================================================
\section{Die Replikationskrise}
%==============================================================================
\nsl{Was geschieht, wenn eine Disziplin diese Gegenmittel zu wenig nutzt,
hat die Psychologie in den 2010er Jahren an sich selbst untersucht. Die
\emph{Open Science Collaboration} versuchte, 100 publizierte Studien aus
führenden psychologischen Zeitschriften zu replizieren
\parencite{opensciencecollaboration2015}. Während 97~Prozent der
Originalstudien signifikante Effekte berichtet hatten, erreichten nur
etwa 36~Prozent der Replikationen ebenfalls Signifikanz, und die
Effektstärken halbierten sich im Mittel. Schon zehn Jahre vorher hatte
John Ioannidis argumentiert, dass bei kleinen Studien, kleinen Effekten,
flexiblen Designs und starkem Interesse an einem bestimmten Ergebnis die
Mehrzahl publizierter Befunde falsch sein kann \parencite{ioannidis2005false}.
Die Replikationskrise ist deshalb weniger eine Krise einzelner Forscher
als ein institutionelles Bias-Problem: Wenn Zeitschriften vor allem
positive, überraschende Ergebnisse drucken und Karrieren an Publikationen
hängen, belohnt das System genau die Verzerrungen, die die Methode
eigentlich ausschalten soll. Die Antwort der Open-Science-Bewegung sind
die Gegenmittel aus dem vorigen Abschnitt, institutionell verankert.}
\begin{description}
\item[{\rot\bf Die Replikationskrise}]~\\[-\lblskip]
\bi
  \ite Open Science Collaboration 2015: 100 psychologische Studien
    repliziert
  \ite Originale: 97\,\% signifikant -- Replikationen: etwa 36\,\%
  \ite Effektstärken im Mittel halbiert
  \ite Ursache: kleine Stichproben, flexible Auswertung, Publikationsdruck
  \ite ein \emph{institutionelles} Bias-Problem -- Antwort: Open Science
\ei
\os{\newpage}
\end{description}

%==============================================================================
\section{Bias-Jagd}
%==============================================================================
\begin{description}
\item[{\rot\bf Aufwärmen (Einzelarbeit, 5 Min.)}]~\\[-\lblskip]
\bi
  \ite Übungsblatt Teil~A, Aufgabe~1: ein fiktiver Kurzbericht
  \ite Markieren Sie jede Stelle, an der ein Bias das Ergebnis
    beeinflusst haben könnte
  \ite Ordnen Sie sie einer Phase zu: Planung, Erhebung, Auswertung,
    Interpretation, Publikation
\ei
\os{\newpage}

\item[{\rot\bf Bias-Jagd im Artikel (Zweiergruppen, 25 Min.)}]~\\[-\lblskip]
\bi
  \ite Übungsblatt Teil~A, Aufgabe~2: Raster
  \ite mindestens \textbf{zwei} Bias-anfällige Stellen finden
  \ite Seite und Zitat angeben, Bias benennen, begründen
  \ite Welches Gegenmittel hätte geholfen? Wurde es verwendet?
  \ite fair bleiben: Was haben die Autorinnen \emph{gut} gemacht?
\ei
\os{\newpage}

\item[{\rot\bf Selbstbezug (Einzelarbeit, 10 Min.)}]~\\[-\lblskip]
\bi
  \ite Übungsblatt Teil~B, Aufgabe~3; Ihre Design-Skizze aus LE~6 vor
    sich legen
  \ite Welcher Bias droht in Ihrem eigenen Forschungsprozess am
    stärksten?
  \ite In welcher Phase?
  \ite Mit welchem \emph{Verfahren} begrenzen Sie ihn?
\ei
\os{\newpage}
\end{description}

\begin{lehrhinweis}
Bei der Bias-Jagd neigen Gruppen dazu, jede Schwäche eines Artikels
\enquote{Bias} zu nennen. Hilfreich ist die Rückfrage, in welche Richtung
die Schwäche das Ergebnis verschieben würde: Eine kleine Stichprobe allein
ist kein Bias, sondern eine Quelle von Unsicherheit; zum Bias wird sie
erst zusammen mit flexibler Auswertung oder selektivem Berichten. Gute
Antworten zitieren die Stelle im Artikel, benennen die Verzerrung,
begründen die Richtung und schlagen ein konkretes Gegenmittel vor. Darauf
achten, dass die Kritik fair bleibt: Ziel ist die Lesekompetenz, nicht
das Vorführen der Autorinnen und Autoren. Beim Selbstbezug ist die
häufigste schwache Antwort \enquote{Ich werde darauf achten,
unvoreingenommen zu sein}; sie wird mit der Leitthese der Einheit
zurückgegeben: Welches \emph{Verfahren}?
\end{lehrhinweis}

%==============================================================================
\section{Lernziele und Aufgaben}
%==============================================================================
\begin{description}
\item[{\rot\bf Nach dieser Einheit können Sie}]~\\[-\lblskip]
\bi
  \ite erklären, warum Verzerrungen auch bei Wissen um sie bestehen
    bleiben
  \ite die häufigsten Verzerrungen den Phasen des Forschungsprozesses
    zuordnen
  \ite Bias und Noise unterscheiden
  \ite Bias-anfällige Stellen in einem empirischen Artikel erkennen und
    begründen
  \ite konkrete Gegenmittel für Ihren eigenen Forschungsprozess benennen
\ei
\os{\newpage}

\item[{\rot\bf Bis zur nächsten Woche}]~\\[-\lblskip]
\bi
  \ite \textbf{Bias-Jagd} (Aufgabe~2) fertigstellen: mindestens zwei
    Stellen, begründet
  \ite \textbf{Reflexion} (Aufgabe~3): Welcher Bias droht in Ihrem
    eigenen Forschungsprozess, und wie begrenzen Sie ihn?
  \ite \textbf{Tagebuch des Staunens:} eine Situation, in der Sie etwas
    wussten und trotzdem anders wahrgenommen haben
\ei
\os{\newpage}
\end{description}

\handoutinput{handouts/LE07}
